\section{Conventions, definitions and the frozen statements}\label{bp:p0}

This part fixes notation and restates, in mathematical form, the five frozen theorems and the
twenty frozen definitions of the formalization (owner approval D11, 27 October 2026). The Lean
declarations in \texttt{KAM/Frozen/} are authoritative; the statements below are their faithful
transcriptions, for use by the proofs in Parts~\ref{bp:p1}--\ref{bp:p4}.

\subsection{Conventions}\label{bp:p0:conventions}

\begin{itemize}
\item $n\ge 2$ is the number of degrees of freedom (ruling Q5). Angles $\theta\in\R^n$, actions
  $I\in\R^n$, phase points $z=(\theta,I)$. Functions are $1$-periodic in each angle, so
  $\T^n=\R^n/\Z^n$ (ruling Q2).
\item On $\R^m$ and $\C^m$ we use the max norm $\abs{v}=\max_i\abs{v_i}$; on $\C^n\times\C^n$ the
  max of the two components. Linear and multilinear maps carry the induced operator norms; for a
  matrix this is $\abs{M}=\max_i\sum_j\abs{M_{ij}}$ (HL19 \S2.1). Sup norms over a set $U$ are
  written $\norm{f}_U=\sup_U\abs{f}$; on a strip, $\norm{f}_\rho=\norm{f}_{T_\rho}$.
\item The complex strip of width $\rho>0$ is $T_\rho=\{\theta\in\C^n:\abs{\operatorname{Im}\theta_i}<\rho\ \forall i\}$
  (HL19 \S2.1); it is invariant under translations by $\Z^n$ (and by $\R^n$).
\item Fourier characters are $e^{2\pi i k\cdot\theta}$, $k\in\Z^n$, with $\hat u_k=\int_{[0,1]^n}u(\theta)e^{-2\pi ik\cdot\theta}\,d\theta$
  and average $\avg{u}=\hat u_0$ (HL19 \S2.1). Weighted coefficient norms
  $\sum_k\abs{\hat u_k}e^{\abs{k}_1 s}$ use a \emph{coefficient radius} $s=2\pi\rho$
  (\texttt{docs/fourier-conventions.md}).
\item Working norms (decision D12). Sup norms on complex domains are the working norms of
  Parts~\ref{bp:p2}--\ref{bp:p4}, as in HL19 and GJLV05. Weighted coefficient norms enter only
  through the cohomological equation and Fourier-coefficient arguments of Part~\ref{bp:p1}.
\end{itemize}

\subsection{Definitions}\label{bp:p0:defs}

\begin{definition}[Diophantine vectors]\label{bp:p0:diophantine}
For $\gamma,\tau\in\R$, a vector $\omega\in\R^n$ is \emph{$(\gamma,\tau)$-Diophantine},
$\omega\in\Dio_{\gamma,\tau}$, if $\gamma/\abs{k}_1^{\tau}\le\abs{k\cdot\omega}$ for every
$k\in\Z^n\setminus\{0\}$, where $\abs{k}_1=\sum_i\abs{k_i}$.
\frozen{KAM.IsDiophantine}
\source{HL19 (3.1), p.~14; Pos01 (1), p.~5.}
\end{definition}

\begin{definition}[Complex neighbourhoods and the fundamental domain]\label{bp:p0:domains}
For $B\subset\R^n$ and $r>0$, $B_r=\{I\in\C^n:\exists J\in B,\ \abs{I-J}<r\}$; the phase domain
of HEAD is $D(s_0,r_0)=T_{s_0}\times B_{r_0}$. The fundamental domain of $\T^n\times B$ is
$F_B=[0,1)^n\times B$, with product Lebesgue measure.
\frozen{KAM.complexStrip, KAM.actionNbhd, KAM.fundamentalDomain}
\end{definition}

\begin{definition}[Hamiltonian vector fields, symplectic form]\label{bp:p0:hamiltonian}
For $H:\R^n\times\R^n\to\R$, $X_H(z)=(\partial_I H(z),-\partial_\theta H(z))$; for holomorphic
$h$ on an open subset of $\C^n\times\C^n$, $X_h$ is defined by the same formula with complex
derivatives. The standard symplectic form is
$\omega_0((a,b),(c,d))=\sum_i(b_ic_i-a_id_i)$, i.e.\ $\omega_0=\sum_i dI_i\wedge d\theta_i$, and
$\Omega_0(a,b)=(-b,a)$ is its matrix $\left(\begin{smallmatrix}0&-I\\ I&0\end{smallmatrix}\right)$.
\frozen{KAM.hamiltonianVectorField, KAM.hamiltonianVectorFieldℂ, KAM.symplecticForm, KAM.omegaZero}
\source{HL19 Rem.~2.1 and \S2.2, p.~6.}
\end{definition}

\begin{remark}[Lean totality]\label{bp:p0:totality}
The frozen definitions use \Lean{fderiv}, which is $0$ where a function is not differentiable,
the matrix inverse, which is $0$ for a singular matrix, and the Bochner integral, which is $0$ for
a non-integrable function. Every frozen theorem guards these by hypotheses (differentiability on
the relevant open set, invertibility, integrability through continuity); the proofs below must
establish the guarded facts before using the defining formulas.
\end{remark}

\begin{definition}[The integrable part, the perturbation, the frequency map]\label{bp:p0:nearlyintegrable}
For $h:\C^n\to\C$ and $f:\C^n\times\C^n\to\C$, the real Hamiltonian is
$H(\theta,I)=\operatorname{Re}\bigl(h(I)+f(\theta,I)\bigr)$ at real points, and the frequency map
is $\nabla h(I)$, the real gradient of $I\mapsto\operatorname{Re}h(I)$ at real $I$.
\frozen{KAM.realHamiltonian, KAM.frequencyMap}
\source{Pos01, Classical KAM Theorem, p.~6.}
\end{definition}

\begin{definition}[Invariant Lagrangian torus]\label{bp:p0:torus}
A map $K:\R^n\to\R^n\times\R^n$ is an \emph{invariant Lagrangian torus of $H$ with frequency
$\omega$} if (i) $K$ is real-analytic; (ii) $K(\theta+k)=K(\theta)+(k,0)$ for $k\in\Z^n$
(primary torus); (iii) $K(\theta')=K(\theta)+(k,0)$ implies $\theta'=\theta+k$; (iv) $DK(\theta)$
is injective; (v) $\omega_0(DK(\theta)u,DK(\theta)v)=0$ for all $u,v$; (vi) for all $\theta$ and
$t$, $\frac{d}{dt}K(\theta+t\omega)=X_H(K(\theta+t\omega))$.
\frozen{KAM.IsInvariantLagrangianTorus}
\source{HL19 (2.1), (2.3), Rem.~2.2--2.3, pp.~6--7; ruling Q3.}
\end{definition}

\begin{definition}[The KAM cover]\label{bp:p0:cover}
Given $h,f,B,r_0,\tau,\delta$, the set $\mathcal{K}(\delta)$ consists of the points $z\in K(\R^n)$
for some $K,\omega,\gamma,I_0$ with: $\gamma>0$; $\omega\in\Dio_{\gamma,\tau}$; $K$ an invariant
Lagrangian torus of $H$ with frequency $\omega$; every action $K(\theta)_2$ within $r_0$ of $B$;
$I_0\in B$ and $\nabla h(I_0)=\omega$; and $\abs{K(\theta)-(\theta,I_0)}\le\delta$ for all $\theta$.
\frozen{KAM.kamCover}
\end{definition}

\begin{definition}[Torus data on strips]\label{bp:p0:stripdata}
Let $K:\C^n\to\C^n\times\C^n$, $h$ holomorphic, $\omega\in\R^n$. $K$ is \emph{primary on $T_\rho$}
if $K(\theta+k)=K(\theta)+(k,0)$ for $\theta\in T_\rho$, $k\in\Z^n$, and \emph{real on reals} if
$K(\R^n)\subset\R^n\times\R^n$. The invariance error is $E(\theta)=X_h(K(\theta))-DK(\theta)\omega$.
With $L=DK$ (a $2n\times n$ matrix) and the bilinear transpose, the Gram matrix is $L^\top L$,
$B=(L^\top L)^{-1}$, the normal frame is $N=\Omega_0LB$, $X_N=DX_h(K)N-DN[\omega]$, the torsion is
$T=N^\top\Omega_0X_N$, and $\avg{T}=\int_{[0,1]^n}T(\theta)\,d\theta$.
\frozen{KAM.IsPrimaryOnStrip, KAM.IsRealOnReals, KAM.invarianceError, KAM.dotℂ, KAM.gramMatrix,
KAM.normalFrame, KAM.torsion, KAM.torsionAverage}
\source{HL19 (2.2), (3.2)--(3.9), pp.~7, 16--17; Rem.~2.3, 3.3 (canonical Case III, $A=0$).}
\end{definition}

\subsection{The frozen theorems}\label{bp:p0:theorems}

\begin{theorem}[HEAD: positive-measure KAM with rate]\label{bp:p0:head}
Let $n\ge2$, $\tau>n-1$, $s_0,r_0>0$, and $B\subset\R^n$ open and bounded with
$\meas\{I\in B:\dist(I,B^c)<\delta\}\le C_B\delta$ for all $\delta>0$. Let $h$ be holomorphic on
$B_{r_0}$, real on real points, with $\abs{I-J}\le L\abs{\nabla h(I)-\nabla h(J)}$ on $\bar B$.
Then there are $C,\varepsilon_0>0$ such that for every $f$ holomorphic on $D(s_0,r_0)$,
$1$-periodic in $\theta$, real on real points, and every $\varepsilon\le\varepsilon_0$ with
$\abs{f}\le\varepsilon$ on $D(s_0,r_0)$:
\[\meas\bigl(F_B\setminus\mathcal{K}(C\sqrt\varepsilon)\bigr)\le C\varepsilon^{1/4}.\]
\frozen{KAM.kam\_positive\_measure}
\source{Pos01, Classical KAM Theorem, p.~6, at the rate of decision D8 (ledger D-MEAS-1).}
\end{theorem}

\begin{theorem}[HEAD-Q: qualitative form]\label{bp:p0:headq}
Same data as Theorem~\ref{bp:p0:head} without the boundary condition. For every $\kappa>0$ there is
$\varepsilon_0>0$ such that, for $f,\varepsilon$ as there with $\varepsilon\le\varepsilon_0$,
$\meas(F_B\setminus\mathcal{K}(\kappa))\le\kappa\meas(F_B)$.
\frozen{KAM.kam\_positive\_measure\_qualitative}
\end{theorem}

\begin{theorem}[G1: a-posteriori KAM at fixed frequency]\label{bp:p0:g1}
Let $n\ge2$, $\tau\ge n-1$, $0<\rho_\infty<\rho$, $M,\sigma_K,\sigma_B,\sigma_T,\mu,\bar\gamma>0$.
There are $C_1,C_2>0$ such that the following holds for all $\mathfrak{B},h,K,\omega,\gamma,e$.
Hypotheses: $\mathfrak{B}\subset\C^n\times\C^n$ open and invariant under $\theta\mapsto\theta+k$;
$h$ holomorphic on $\mathfrak{B}$, $1$-periodic in $\theta$, real on real points, with
$\abs{Dh},\abs{X_h},\abs{DX_h},\abs{D^2X_h}\le M$ on $\mathfrak{B}$; $K$ holomorphic on $T_\rho$,
primary, real on reals, with $B_\mu(K(\theta))\subset\mathfrak{B}$ and $\abs{DK(\theta)}\le\sigma_K-\mu$
on $T_\rho$; $L^\top L$ invertible with $\abs{B}\le\sigma_B-\mu$ on $T_\rho$; $\avg{T}$ invertible
with $\abs{\avg{T}^{-1}}\le\sigma_T-\mu$; $0<\gamma\le\bar\gamma$ and $\omega\in\Dio_{\gamma,\tau}$;
$\abs{E}\le e$ on $T_\rho$ and $C_1e/(\gamma^4\rho^{4\tau})<1$.
Conclusion: there is $K_\infty$, holomorphic on $T_{\rho_\infty}$, primary, real on reals, with
$E_\infty=0$ on $T_{\rho_\infty}$, $B_{\mu/2}(K_\infty(\theta))\subset\mathfrak{B}$,
$\abs{DK_\infty}\le\sigma_K-\mu/2$, and $\abs{K_\infty-K}\le C_2e/(\gamma^2\rho^{2\tau})$ on
$T_{\rho_\infty}$.
\frozen{KAM.aposteriori\_kam}
\source{HL19 Theorem 3.1, pp.~16--18 (ledger D-G1-1).}
\end{theorem}

\begin{theorem}[G2: local uniqueness for flows]\label{bp:p0:g2}
Let $n\ge2$, $\tau\ge n-1$, $\rho>0$, $M,\sigma_K,\sigma_B,\sigma_T,\mu>0$. There is $c>0$ such that:
if $\mathfrak{B},h$ are as in G1 (open, $\Z^n$-invariant domain; $h$ holomorphic, periodic, real, with
$\abs{Dh},\abs{X_h},\abs{DX_h},\abs{D^2X_h}\le M$), $\omega\in\Dio_{\gamma,\tau}$ with $\gamma>0$ (no upper
bound on $\gamma$ and no smallness condition), $K_1$ is holomorphic on $T_\rho$, primary, real on reals, with
$B_\mu(K_1(\theta))\subset\mathfrak{B}$, $\abs{DK_1}\le\sigma_K-\mu$, $L^\top L$ invertible with
$\abs{B}\le\sigma_B-\mu$ on $T_\rho$, $\avg{T}$ invertible with $\abs{\avg{T}^{-1}}\le\sigma_T-\mu$, and
$E_1=0$ on $T_\rho$; $K_2$ is holomorphic on $T_\rho$,
primary, real on reals, $K_2(T_\rho)\subset\mathfrak{B}$, $E_2=0$, and
$\abs{K_1-K_2}\le c\gamma^2\rho^{2\tau}$ on $T_\rho$, then $K_2(\theta)=K_1(\theta+\varphi)$ on
$T_{\rho/2}$ for some $\varphi\in\R^n$.
\frozen{KAM.local\_uniqueness\_flow}
\source{Our own result (ledger D-UNIQ-1); template GJLV05 Theorem 2, p.~862.}
\end{theorem}

\begin{theorem}[G3: measure of the non-Diophantine set]\label{bp:p0:g3}
Let $\tau>n-1$ and $\Omega\subset\R^n$ bounded. There is $C$ such that
$\meas\{\omega\in\Omega:\omega\notin\Dio_{\alpha,\tau}\}\le C\alpha$ for every $\alpha>0$.
\frozen{KAM.volume\_not\_diophantine\_le}
\source{Pos01 \S1d, p.~5.}
\end{theorem}
